% Intended LaTeX compiler: lualatex
\documentclass[a4paper, 11pt]{article}

\usepackage{amsmath}
\usepackage{fontspec}
\usepackage{graphicx}
\usepackage{longtable}
\usepackage{wrapfig}
\usepackage{rotating}
\usepackage[normalem]{ulem}
\usepackage{capt-of}
\usepackage{hyperref}
\usepackage[danish]{babel}
\usepackage{mathtools}
\usepackage[margin=2.0cm]{geometry}
\hypersetup{colorlinks, linkcolor=black, urlcolor=blue}
\usepackage{fancyhdr}
\fancyhf{}
\pagestyle{fancy}
\fancyhead[L]{Bevisstafet - 2.gradsligningen} \fancyhead[C]{Matematik A}\fancyhead[R]{Jacob Debel}
\setlength{\parindent}{0em}
\parskip 1.5ex
\date{}
\title{}
\hypersetup{
 pdfauthor={},
 pdftitle={},
 pdfkeywords={},
 pdfsubject={},
 pdfcreator={},
 pdflang={Danish}}
\begin{document}

\section*{Formål}
\label{sec:org13f370d}
Bevisstafetten er en samarbejdsøvelse, hvor I som hold skal udlede løsningsformlen for en andengradsligning:
\begin{align*}
ax^2 + bx + c &= 0 \iff \\
\dots \\
\dots \\
x &= \frac{-b \pm \sqrt{b^2-4ac}}{2a}
\end{align*}

I skal kombinere jeres egen matematiske viden med målrettede "stafetløb" til et udledningsark, når I får brug for hjælp.
\section*{Regler for Stafetten}
\label{sec:org4cae23c}

\subsection*{1. Start selv}
\label{sec:org49335bf}
\begin{itemize}
\item I har \textbf{ét fælles papir} og \textbf{én pen} ved bordet.
\item I begynder blot at udlede beviset så langt, I kan på egen hånd.
\item Hvert hold har en \textbf{løber} og en \textbf{skriver}.
\end{itemize}
\subsection*{2. Skriv kun hvad der dikteres}
\label{sec:org8bb31c5}
\begin{itemize}
\item \textbf{Skriveren} må \textbf{kun} skrive det ned på papiret, som de andre på holdet forklarer og dikterer.
\item \textbf{Skriveren} må ikke selv gætte eller regne videre uden om holdet.
\end{itemize}
\subsection*{3. Hent hjælp ved behov}
\label{sec:org973aafa}
\begin{itemize}
\item Hvis I går i stå eller bliver i tvivl, sender I \textbf{løberen} op til udledningsarket.
\item \textbf{Ingen noter:} Løberen må \textbf{ikke} medbringe papir, mobil eller pen op til udledningsarket. Alt skal huskes i hovedet.
\end{itemize}
\subsection*{4. Tavshed indtil bordet}
\label{sec:org89ee044}
\begin{itemize}
\item Løberen må \textbf{først sige noget}, når vedkommende er helt tilbage ved holdets bord.
\end{itemize}
\subsection*{5. Skift på tur}
\label{sec:orgfaa143c}
\begin{itemize}
\item Rollerne som \textbf{løber} og \textbf{skriver} skal gå på tur i gruppen, hver gang I henter ny information.
\end{itemize}
\section*{Vinderbetingelse}
\label{sec:org8d414f4}
Når I mener, at I har udledt hele formlen korrekt og med forklarende tekstbegrundelser, kalder I på læreren. Læreren laver en hurtig stikprøve hos en tilfældig elev på holdet, som skal kunne forklare et udvalgt trin.
\end{document}
